% Portuguese version of Curriculum_Vitae_English.tex — keep entries in sync with it.
\documentclass[a4paper,9pt]{article}
\usepackage[left=1.5cm,right=1.5cm,top=2cm,bottom=2cm]{geometry}
\usepackage{fancyhdr}
\usepackage{tabularx}
\usepackage[utf8]{inputenc}
\usepackage[brazil]{babel}
\usepackage{verbatim}
\usepackage[T1]{fontenc}
\usepackage{lmodern}            % scalable fonts (needed by microtype expansion)
\usepackage{CV}
\usepackage{amstext}
\usepackage{enumitem}
\usepackage{microtype}          % subtle spacing/justification improvements
\usepackage{hyperref}
\usepackage{fontawesome5}
\usepackage{marvosym}
\usepackage{academicons}

\hypersetup{breaklinks=true, pagecolor=white, colorlinks=true, citecolor=blue, linkcolor=black, urlcolor=black}
\urlstyle{rm} % suppresses typewriter font for url

% Tighter, consistent lists throughout
\setlist[itemize]{leftmargin=1.5em, itemsep=2pt, topsep=2pt, parsep=0pt}

% Uniform spacing between publication / working-paper entries
\newcommand{\cvgap}{\vspace{2mm}}

\fancyhead[L]{\small{\emph{Curriculum Vitae}}}
\fancyhead[C]{\Large{\textsc{Francisco Cavalcanti}}}
\fancyhead[R]{\small{\emph{\today}}}
\pagestyle{fancy}
\setlength{\headheight}{18pt}
\begin{document}
\vspace*{0.1mm}
\normalsize

\section{Informações Pessoais}

\begin{flushleft}
\begin{tabular}{lll}
\faGlobe\ \href{https://franciscocavalcanti.github.io}{franciscocavalcanti.github.io}   &\hspace{2mm} \Letter\hspace{0.2mm} \href{mailto:francisco.dlcavalcanti@ufpe.br}{\nolinkurl{francisco.dlcavalcanti@ufpe.br}}  &\hspace{2mm} \\
\end{tabular}
\end{flushleft}

\section{Cargos}

\begin{flushleft}
\begin{tabular}{ll}
\textbf{Professor Adjunto do \href{https://www.ufpe.br/decon}{Departamento de Economia da UFPE} (2023 -- presente)} \\
Pesquisador de pós-doutorado no \href{http://www.econ.puc-rio.br/}{Departamento de Economia da PUC-Rio} (2018 -- 2023) \\
Coordenador operacional do \href{http://www.econ.puc-rio.br/datazoom/index.html}{Data Zoom} (2019 -- 2023) \\
\end{tabular}
\end{flushleft}

\section{Formação Acadêmica}

\begin{itemize}

\item{\textbf{Doutorado} \emph{excellent cum laude} em Economia pela \href{http://www.ub.edu/school-economics/}{Universidade de Barcelona} (2018). \\ Doutorando visitante na \href{https://warwick.ac.uk/fac/soc/economics/}{Universidade de Warwick} (Reino Unido) -- segundo semestre de 2016. \\ Orientadores: Amedeo Piolatto, Albert Sol\'e-Oll\'e.}

\item{\textbf{Mestrado} em Economia pelo \href{https://sites.google.com/view/pimes/principal}{PIMES/UFPE} (2013). Orientador: Raul Silveira-Neto.}

\item{\textbf{Bacharelado} em Economia pela \href{https://www.ufpe.br/ccsa/}{Universidade Federal de Pernambuco} (2010).}

\end{itemize}

\section{Publicações Selecionadas}

\textbf{\href{https://doi.org/10.1016/j.jebo.2025.106998}{Natural Disasters and Voting Behavior Under Authoritarian Regimes: Evidence From the Brazilian Shrimp Vote.}} \emph{Journal of Economic Behavior \& Organization, 2025}, Vol.~234, p.~106998. (com D.~Baerlocher, R.~Caldas \& R.~Schneider)

\cvgap

\textbf{\href{https://onlinelibrary.wiley.com/doi/full/10.1111/obes.12560}{Drought-Reliefs and Partisanship.}} \emph{Oxford Bulletin of Economics and Statistics, 2024}, Vol.~86, p.~187--208. (com \href{https://sites.google.com/site/federicoboffaunibz/}{F.~Boffa}, \href{https://sites.google.com/site/piolatto/home}{A.~Piolatto} \& \href{https://sites.google.com/site/cfonsrosen/home}{C.~Fons-Rosen}) \\
\hspace{2em} \href{https://voxeu.org/article/insights-distributive-politics-new-way-measure-aridity}{VoxEU: New insights on distributive politics from a new way to measure aridity} \\
\hspace{2em} \href{https://github.com/FranciscoCavalcanti/Drought-reliefs-and-Partisanship}{Repositório no GitHub}

\cvgap

\textbf{\href{https://www.sciencedirect.com/science/article/pii/S0047272718301361}{Popularity Shocks and Political Selection.}} \emph{Journal of Public Economics, 2018}, Vol.~165, p.~201--216. (com \href{http://www.gianmarcodaniele.com/}{G.~Daniele} \& \href{http://sergio-galletta.com/}{S.~Galletta})

\section{Capítulos de Livro}

\textbf{\href{https://doi.org/10.18235/0014396}{Total Factor Productivity Growth in Brazilian Agriculture (1985--2017): The Roles of Climate Change and Public Policy.}} In: \emph{Agricultural Productivity in Latin America and the Caribbean: What We Know and Where We Are Heading}, organizado por L.~Salazar, M.~Schling, C.~P.~De Salvo \& P.~Martel. Banco Interamericano de Desenvolvimento, 2026. (com S.~Helfand, C.~de Freitas, A.~Moreira \& M.~Schling)

\section{Textos para Discussão (Working Papers)}

\textbf{\href{https://www.researchgate.net/profile/Francisco-Cavalcanti-6/publication/406502168_Education_and_Job_Loss_During_COVID-19_Evidence_from_Brazil/links/6a29d64a68552924bb0049a2/Education-and-Job-Loss-During-COVID-19-Evidence-from-Brazil.pdf}{Education and Job Loss During COVID-19: Evidence from Brazil.}} (com F.~Didier \& G.~Gonzaga)

\cvgap

\textbf{\href{https://ideas.repec.org/p/usf/wpaper/2026-01.html}{Beyond Appearance: The Socioeconomic and Historical Roots of Racial Identity in Brazil.}} \emph{University of South Florida, Department of Economics Working Papers, 2026.} (com D.~Baerlocher \& R.~Caldas)

\cvgap

\textbf{\href{https://www.researchgate.net/profile/Ricardo-Carvalho-10/publication/383129940_Do_Landslides_Impact_Urbanization_Patterns_Evidence_from_Brazilian_Cities/links/66be0c688d0073559255b1e4/Do-Landslides-Impact-Urbanization-Patterns-Evidence-from-Brazilian-Cities.pdf}{Do Landslides Impact Urbanization Patterns? Evidence from Brazilian Cities.}} (2024, com R.~Lima, P.~Alves \& B.~Barsanetti)

\cvgap

\textbf{\href{https://economics.ucr.edu/docs/helfand/Helfand\%20Climate\%20Change,\%20Drought,\%20and\%20Ag\%20Prod\%20in\%20Brazil\%208-2024.pdf}{Climate Change, Drought, and Agricultural Production in Brazil.}} (com S.~Helfand \& A.~Moreira)

\cvgap

\textbf{\href{http://diposit.ub.edu/dspace/handle/2445/182602}{Ignorance is Bliss: Voter Education and Alignment in Distributive Politics.}} \emph{IEB Working Paper 2021/07, Institut d'Economia de Barcelona, 2021.} (com F.~Boffa \& A.~Piolatto)

\cvgap

\textbf{\href{https://mpra.ub.uni-muenchen.de/88317/}{Voters Sometimes Provide the Wrong Incentives: The Lesson of the Brazilian Drought Industry.}} (2018)

\cvgap

\textbf{\href{https://www.researchgate.net/publication/319902513_Creative_Class_Human_Capital_and_Urban_Dynamism_Empirical_Evidence_for_the_Brazilian_Cities}{Creative Class, Human Capital and Urban Dynamism: Empirical Evidence for the Brazilian Cities.}} \emph{Anais do XLII Encontro Nacional de Economia, n.~160, 2016.} (com R.~Silveira-Neto)

\section{Trabalhos em Andamento}

\begin{itemize}
\item Droughts, Sanitation Provision, and Access to Services in Brazil (com T.~Soares \& P.~F.~Azevedo)
\item Mobile Broadband Expansion and Tasks: Evidence from Brazilian Formal Labor Markets (com H.~Mota \& G.~Gonzaga)
\item Drought and Corruption
\item Show Me Your Pix: The Impact of Electronic Payment Technologies on Tax Collection
\end{itemize}

\section{Auxílios e Bolsas}

\noindent \textbf{PROPESQI/UFPE} -- Apoio a Jovens Pesquisadores (2026). \textbf{ESRC} UK-Brazil global talent exchange scheme (2025). \textbf{CAPES} -- bolsista de pós-doutorado (2018 -- 2023). \textbf{AUXÍLIO} -- Ministério da Economia e Competitividade / Governo da Espanha (ECO2015-68311-R). Pesquisador principal: Albert Sol\'e-Oll\'e (2016 -- 2018). \textbf{APIF} -- Bolsa de doutorado da UB (2015 -- 2018). \textbf{Bolsa de doutorado do IEB} -- Institut d'Economia de Barcelona (2014 -- 2015). \textbf{AUXÍLIO} -- Ministério da Economia e Competitividade / Governo da Espanha (ECO2012-37131). Pesquisador principal: Albert Sol\'e-Oll\'e (2013 -- 2015). \textbf{Bolsa CNPq} -- Ministério da Educação do Brasil (2013 -- 2013). \textbf{Bolsa de mestrado CNPq} -- Ministério da Educação do Brasil (2011 -- 2013).

\vspace*{0.1mm}

\section{Conferências, Workshops e Seminários}

\noindent \textbf{2026} CAF -- Workshop do Banco de Desenvolvimento da América Latina e do Caribe (Buenos Aires, Argentina), IIPF -- Congresso Anual do International Institute of Public Finance (Lisboa, Portugal, \emph{previsto}), MadBar Workshop on Public Economics (Barcelona, Espanha, \emph{previsto}). \textbf{2025} Just Transition Hub (Dundee, Reino Unido), Liverpool School of Tropical Medicine (Liverpool, Reino Unido), Universidade de York (York, Reino Unido). \textbf{2024} Society for Institutional \& Organizational Economics (SIOE) e Ronald Coase Institute Alumni Conference (Chicago, EUA), Seminários de Economia do PIMES/UFPE (Recife, Brasil). \textbf{2023} Seminários de Economia do PIMES/UFPE (Recife, Brasil), Seminário de Economia do PPGE/UFPB (João Pessoa, Brasil). \textbf{2022} LACEA-LAMES (Lima, Peru). \textbf{2021} 2º Workshop AMZ 2030 (Belém, Brasil). \textbf{2020} Seminários de Economia da FEA-USP (Online), Seminários de Economia da UCB (Online). \textbf{2019} Seminários de Economia do PIMES/UFPE (Recife, Brasil), Ronald Coase Workshop on Institutional Analysis (Varsóvia, Polônia). \textbf{2018} Higher School of Economics (HSE) (Moscou, Rússia), Carioca Workshop on Political Economy (Rio de Janeiro, Brasil), SBE -- Sociedade Brasileira de Econometria (Rio de Janeiro, Brasil). \textbf{2017} SAEe -- Associação Espanhola de Economia (Barcelona, Espanha), Seminários de Economia do PIMES/UFPE (Recife, Brasil), LACEA-LAMES (Buenos Aires, Argentina). \textbf{2016} Seminário CWIP / Universidade de Warwick (Coventry, Reino Unido), Workshop on The Political Economy of Federalism and Local Development (Bruneck, Itália), Conference on Development Economics and Policy (Heidelberg, Alemanha), Seminários de Economia do PIMES/UFPE (Recife, Brasil), Seminários do IEB (Barcelona, Espanha), Lisbon Meeting on Institutions and Political Economy (Lisboa, Portugal). \textbf{2015} Workshop on PhD in Economics of UB (Barcelona, Espanha), Summer School in Development Economics (Garda, Itália). \textbf{2014} VII Italian Doctoral Workshop in Empirical Economics (Turim, Itália), Seminários de Economia do PIMES/UFPE (Recife, Brasil), Seminários do IEB (Barcelona, Espanha).

\section{Atividades Acadêmicas e de Pesquisa}

\begin{flushleft}
\begin{tabular}{l}
\textbf{Parecerista (Referee)}: Economia, British Journal of Political Science, The Economic Journal, \\
The Review of Economics and Statistics \\
\\
\textbf{Comissão Organizadora}: Encontro Anual da LACEA 2025, Seminário de Microeconomia \\
Aplicada da PUC-Rio (2019 -- 2020), Seminário de Doutorado em Economia da \\
Universidade de Barcelona (2014 -- 2016). \\
\end{tabular}
\end{flushleft}

\section{Ensino}

\begin{flushleft}
\begin{tabular}{lll}
\multicolumn{3}{l}{\textbf{Pós-graduação}} \\
Seminários de Discentes & 2025 -- 2026 & UFPE \\
\\
\multicolumn{3}{l}{\textbf{Graduação}} \\
Macroeconomia & 2026 & UFPE \\
Economia e Mudança Climática & 2025 -- 2026 & UFPE \\
Economia 12 & 2024 -- 2026 & UFPE \\
Economia Monetária & 2024 -- 2025 & UFPE \\
Microeconomia & 2024 & UFPE \\
Introdução à Economia & 2023 & UFPE \\
Economia Política & 2017 & Universidade de Barcelona \\
Fundamentos de Tributação & 2015 -- 2017 & Universidade de Barcelona \\
Macroeconomia & 2013 & UFPE \\
Macroeconomia & 2012 & UFRPE \\
\end{tabular}
\end{flushleft}

\vspace*{0.1mm}

\section*{Relatórios Técnicos e Outras Publicações}

\begin{flushleft}

\textbf{Análise das Desigualdades Socioeconômicas de Gênero e Raça e Políticas Públicas na Amazônia Brasileira.} Consultoria -- Proyecto RG-T3975: Desarrollo Sostenible de la Amazonía con Enfoque de Género y Diversidad, setembro de 2025.

\cvgap

\textbf{\href{https://amazonia2030.org.br/desigualdades-no-mercado-de-trabalho-por-raca-evidencias-para-a-amazonia-legal/}{Desigualdades no Mercado de Trabalho por Raça: Evidências para a Amazônia Legal.}} (com Gustavo Gonzaga). Projeto Amazônia 2030, n.~42, junho de 2022, p.~42.

\cvgap

\textbf{\href{https://amazonia2030.org.br/desigualdades-no-mercado-de-trabalho-por-genero-evidencias-para-a-amazonia-legal/}{Desigualdades no Mercado de Trabalho por Gênero: Evidências para a Amazônia Legal.}} (com Gustavo Gonzaga). Projeto Amazônia 2030, n.~36, abril de 2022, p.~42.

\cvgap

\textbf{\href{https://amazonia2030.org.br/dinamismo-de-emprego-e-renda-na-amazonia-legal-ocupacoes-qualificadas-e-de-lideranca/}{Dinamismo de Emprego e Renda na Amazônia Legal: Ocupações Qualificadas e de Liderança.}} (com Flávia Alfenas \& Gustavo Gonzaga). Projeto Amazônia 2030, n.~15, agosto de 2021.

\cvgap

\textbf{\href{https://amazonia2030.org.br/dinamismo-de-emprego-e-renda-na-amazonia-legal-servicos/}{Dinamismo de Emprego e Renda na Amazônia Legal: Setor Serviços.}} (com Flávia Alfenas \& Gustavo Gonzaga). Projeto Amazônia 2030, n.~14, agosto de 2021.

\cvgap

\textbf{\href{https://amazonia2030.org.br/dinamismo-de-emprego-e-renda-na-amazonia-legal-setor-publico/}{Dinamismo de Emprego e Renda na Amazônia Legal: Setor Público.}} (com Flávia Alfenas \& Gustavo Gonzaga). Projeto Amazônia 2030, n.~13, agosto de 2021.

\cvgap

\textbf{\href{https://amazonia2030.org.br/dinamismo-de-emprego-e-renda-na-amazonia-legal-agropecuaria/}{Dinamismo de Emprego e Renda na Amazônia Legal: Agropecuária.}} (com Flávia Alfenas \& Gustavo Gonzaga). Projeto Amazônia 2030, n.~10, agosto de 2021.

\cvgap

\textbf{\href{https://amazonia2030.org.br/mercado-de-trabalho-na-amazonia-legal-uma-analise-comparativa-com-o-resto-do-brasil/}{Mercado de Trabalho na Amazônia Legal: Uma Análise Comparativa com o Resto do Brasil.}} (com Flávia Alfenas \& Gustavo Gonzaga). Projeto Amazônia 2030, novembro de 2020, p.~68.

\cvgap

\textbf{Nível e Evolução da Desigualdade de Renda na Bahia: uma avaliação do papel da educação e dos programas sociais.} (com Rodrigo Oliveira). Bahia Análise \& Dados, v.~24, p.~89, 2014.

\end{flushleft}

\end{document}
